% !TEX root =./ECE448_Practice_main.tex
%
% L4 practice set (SDR Introduction). Blank copy: \begin{solution}[h] reserves
% h of work space. SOLUTIONS copy (\iskey=1) prints the red worked solutions.
% Radio specifications are the manufacturers' numbers, as on the lesson page.

\fullwidth{\textbf{Learning Objective 2.1: } I can describe what a Software Defined
Radio is, contrast it with a traditional hardware radio, and identify the SDR platforms
used in this course (RTL-SDR, HackRF, PlutoSDR, ADALM-PHASER).
\\[2pt]\textbf{Learning Objective 2.2: } I can identify the roles of the RF frontend,
ADC/DAC, and host in an SDR signal chain and describe sample-rate, bandwidth, and
dynamic-range trade-offs.
\\[4pt]\normalfont\small Use the course SDR table from the lesson page. Show your work.
Write a \textbf{Documentation} line at the end (\textbf{None} if you did not
collaborate).}

\begin{questions}

%%%%%%%%%%%%%%%%%%%%%%%%%%%%%%%%%%%%%%%%%
\question \textbf{Choosing a radio.} The course SDRs are the RTL-SDR Blog V4
(0.5 to 1766 MHz, 2.56 MS/s, 8 bits, receive only), the HackRF One (1 to 6000 MHz,
20 MS/s, 8 bits, half duplex), the PlutoSDR (325 to 3800 MHz, about 4 MS/s sustained over USB,
12 bits, full duplex), and the ADALM-PHASER (10.0 to 10.5 GHz, digitized by a PlutoSDR).
\begin{parts}

	\part For each task, name a course SDR that can do it, and give the specification
	that rules out each of the others.
	\begin{enumerate}[label=(\roman*)]
		\item Receive ship AIS at 162 MHz.
		\item Display the whole FM broadcast band, 88 to 108 MHz, at once.
		\item Transmit a 915 MHz test signal into a cable and receive it at the same time.
		\item Find the direction of a 10.3 GHz beacon.
	\end{enumerate}
		\begin{solution}[2.3in]
		(i) The RTL-SDR or the HackRF; 162 MHz is inside both tuning ranges. The
		PlutoSDR starts at 325 MHz and the ADALM-PHASER works only at 10.0 to 10.5 GHz.

		(ii) The HackRF at 20 MS/s, which shows a 20 MHz window. That leaves no
		margin, so the stations at the very edges of the band sit in the filters'
		roll-off. The RTL-SDR's 2.56 MS/s shows only 2.56 MHz, the PlutoSDR cannot tune
		to 100 MHz, and the ADALM-PHASER works only at X band.

		(iii) The PlutoSDR, the only full-duplex radio. The HackRF transmits or receives
		but not both at once, the RTL-SDR cannot transmit, and the ADALM-PHASER does
		not cover 915 MHz.

		(iv) The ADALM-PHASER: 10.3 GHz is inside its band, its 8 elements let it find
		direction, and every other radio stops at 6 GHz or below.
		\end{solution}

	\part Explain how an SDR differs from a hardware FM radio, and why the same
	RTL-SDR can decode both broadcast FM and aircraft ADS-B.
		\begin{solution}[1.1in]
		A hardware radio builds its filters and demodulator as circuits for one kind of
		signal. An SDR only tunes, filters, and digitizes a band, and software does the
		filtering, demodulation, and decoding. Both 100 MHz and 1090 MHz are inside the
		RTL-SDR's tuning range, so receiving either one is a matter of setting the
		center frequency and running a different flowgraph.
		\end{solution}

\end{parts}

\newpage
%%%%%%%%%%%%%%%%%%%%%%%%%%%%%%%%%%%%%%%%%
\question \textbf{The mixer and I and Q.} We set a PlutoSDR's center frequency, its LO,
to 433.0 MHz, and a key fob transmits at 433.92 MHz.
\begin{parts}

	\begin{minipage}[t]{0.58\linewidth}
		\part Find the two frequencies at the mixer output, and say which one the
		low-pass filter keeps.
		\begin{solution}[1.0in]
		\eq{ 433.92 - 433.0 &= 0.92\ \text{MHz} \\ 433.92 + 433.0 &= 866.92\ \text{MHz} }
		The low-pass filter keeps the 0.92 MHz difference and removes the sum.
		\end{solution}
	\end{minipage}\hfill%
	\begin{minipage}[t]{0.37\linewidth}
		\raggedleft \vspace{12pt}
		kept: \ansbox[1.2in]{0.92 MHz}

		\vspace{4pt}
		removed: \ansbox[1.2in]{866.92 MHz}
	\end{minipage}

	\begin{minipage}[t]{0.58\linewidth}
		\part With a single real mixer, a second transmitter would also come out at
		0.92 MHz. Find its frequency, and explain how the PlutoSDR's $I$ and $Q$ outputs
		tell the two apart.
		\begin{solution}[1.3in]
		\eq{ 433.0 - 0.92 = 432.08\ \text{MHz} }
		The second mixer, driven by the LO shifted 90°, gives $Q$. In $I + jQ$ the
		433.92 MHz signal rotates one way and appears at $+0.92$ MHz, and the
		432.08 MHz signal rotates the other way and appears at $-0.92$ MHz.
		\end{solution}
	\end{minipage}\hfill%
	\begin{minipage}[t]{0.37\linewidth}
		\raggedleft \vspace{12pt}
		\ansbox[1.2in]{432.08 MHz}
	\end{minipage}

\end{parts}

\newpage
%%%%%%%%%%%%%%%%%%%%%%%%%%%%%%%%%%%%%%%%%
\question \textbf{Sample rate and data rate.} A team captures a 10 MHz-wide test signal
centered on 2437 MHz with a HackRF. The usable window is somewhat narrower than $f_s$,
because the radio's filters roll off near its edges.
\begin{parts}

	\begin{minipage}[t]{0.58\linewidth}
		\part Find the lowest sample rate whose window just holds the signal. Then
		choose a rate 25\% higher, to leave margin at the edges, and find the band
		that window covers.
		\begin{solution}[1.1in]
		A complex window is $f_s$ wide, so the signal needs at least 10 MS/s. With
		25\% margin:
		\eq{ f_s &= 1.25 \times 10 = 12.5\ \text{MS/s} \\ 2437 \pm 6.25 &= 2430.75 \text{ to } 2443.25\ \text{MHz} }
		\end{solution}
	\end{minipage}\hfill%
	\begin{minipage}[t]{0.37\linewidth}
		\raggedleft \vspace{12pt}
		lowest: \ansbox[1.3in]{10 MS/s}

		\vspace{4pt}
		chosen: \ansbox[1.3in]{12.5 MS/s}

		\vspace{4pt}
		band: \ansbox[1.6in]{2430.75 to 2443.25 MHz}
	\end{minipage}

	\begin{minipage}[t]{0.58\linewidth}
		\part GNU Radio stores each complex sample in 8 bytes. At 12.5 MS/s, find the
		data rate through the flowgraph and the size of a 10 s File Sink capture.
		\begin{solution}[0.9in]
		\eq{ 8 \times 12.5\times10^{6} &= 100\ \text{MB/s} \\ 10 \times 100 &= 1\ \text{GB} }
		\end{solution}
	\end{minipage}\hfill%
	\begin{minipage}[t]{0.37\linewidth}
		\raggedleft \vspace{12pt}
		rate: \ansbox[1.2in]{100 MB/s}

		\vspace{4pt}
		file: \ansbox[1.2in]{1 GB}
	\end{minipage}

	\begin{minipage}[t]{0.58\linewidth}
		\part The HackRF sends 8-bit $I$ and 8-bit $Q$ over USB. Find the USB data rate
		at 12.5 MS/s, and how many times larger the flowgraph's rate is.
		\begin{solution}[0.9in]
		\eq{ 2\ \text{bytes} \times 12.5\times10^{6} = 25\ \text{MB/s}, \qquad \frac{100}{25} = 4 }
		\end{solution}
	\end{minipage}\hfill%
	\begin{minipage}[t]{0.37\linewidth}
		\raggedleft \vspace{12pt}
		USB: \ansbox[1.2in]{25 MB/s}

		\vspace{4pt}
		ratio: \ansbox[1.2in]{4}
	\end{minipage}

	\part A teammate suggests the RTL-SDR instead. Give two separate reasons it
	cannot make this capture.
		\begin{solution}[0.9in]
		Its tuning range stops at 1766 MHz, below 2437 MHz. Even inside its range, its
		2.56 MS/s shows only 2.56 MHz, far less than the 10 MHz signal.
		\end{solution}

\end{parts}

\newpage
%%%%%%%%%%%%%%%%%%%%%%%%%%%%%%%%%%%%%%%%%
\question \textbf{Bits and clipping.} Use the L4 flowgraph, \texttt{L04-adc-bits.grc},
and its plots on the lesson page.
\begin{parts}

	\begin{minipage}[t]{0.58\linewidth}
		\part How many levels does the ADC block round to with 8 bits, the RTL-SDR's
		and HackRF's resolution, and with 12 bits, the PlutoSDR's?
		\begin{solution}[0.8in]
		\eq{ 2^{8} = 256, \qquad 2^{12} = 4096 }
		\end{solution}
	\end{minipage}\hfill%
	\begin{minipage}[t]{0.37\linewidth}
		\raggedleft \vspace{12pt}
		8 bits: \ansbox[1.2in]{256}

		\vspace{4pt}
		12 bits: \ansbox[1.2in]{4096}
	\end{minipage}

	\part At 8 bits the weak tone at $-0.49$ MHz is not visible. Switching the menu to
	12 bits makes it appear, about 40 dB above the floor. Did the weak signal change?
	Describe what changed on the spectrum, and why.
		\begin{solution}[1.1in]
		The weak signal did not change; the input is the same. The floor dropped by about
		20 dB and the spurs disappeared, because 12 bits round in much finer steps, so the
		rounding error is smaller. The weak tone was always there, below the 8-bit floor.
		\end{solution}

	\part On a real 8-bit radio, a strong local station sits at $-10$ dBFS and a weak
	signal we want is hidden in the floor. How far can the gain rise before the strong
	station clips? Would raising it that far reveal the weak signal? Explain, and name
	one thing that would help.
		\begin{solution}[1.3in]
		10 dB, since 0 dBFS is full scale. Gain raises the strong and the weak signal
		together, so the weak one rises at most 10 dB before the strong one clips. If
		it sits deeper than that in the floor, it stays hidden. A radio with more bits
		helps, or a narrower window or a filter that keeps the strong station out of
		the ADC.
		\end{solution}

	\part In the L4 flowgraph, name two things on the plots that show the ADC is
	clipping.
		\begin{solution}[0.8in]
		The time sink shows $I$ and $Q$ flattened at $\pm 1$, and the spectrum shows new
		signals where nothing transmits. The strong tone also rises by less than
		the gain increase.
		\end{solution}

\end{parts}

\vspace{1.5em}
\noindent\textbf{Documentation:} \rule[-2pt]{0.6\linewidth}{0.4pt}

\end{questions}
